\documentclass[11pt]{article}
\usepackage[margin=1in]{geometry}
\usepackage{parskip}
\usepackage[hidelinks]{hyperref}
\usepackage{enumitem}
\usepackage{tabularx}
\usepackage{array}
\usepackage{titlesec}
\usepackage{lmodern}
\usepackage[T1]{fontenc}

\titleformat{\section}{\large\bfseries}{\thesection.}{0.5em}{}
\titleformat{\subsection}{\normalsize\bfseries}{\thesubsection.}{0.5em}{}

\title{Authoritative Sources for Patterns of Enterprise Agentic Application Architecture}
\author{}
\date{}

\begin{document}
\maketitle

\section{Purpose}
This handout compiles authoritative sources for studying \textit{Patterns of Enterprise Agentic Application Architecture}. It emphasizes enterprise-grade references that cover agentic design patterns, architecture layers, orchestration approaches, and the historical context of enterprise application architecture.

\section{Primary Sources}

\begin{tabularx}{\textwidth}{>{\bfseries}p{0.22\textwidth}X}
Salesforce & \href{https://architect.salesforce.com/docs/architect/fundamentals/guide/enterprise-agentic-architecture.html}{Enterprise Agentic Architecture and Design Patterns}. A detailed enterprise-focused pattern library for agentic systems, including interaction, specialist, utility, and long-running patterns, plus orchestration archetypes and governance guidance. \\
Google Cloud & \href{https://docs.cloud.google.com/architecture/choose-design-pattern-agentic-ai-system}{Choose a design pattern for your agentic AI system}. A structured architecture guide covering single-agent, multi-agent, sequential, parallel, loop, coordinator, hierarchical decomposition, swarm, ReAct, and human-in-the-loop patterns. \\
Martin Fowler & \href{https://martinfowler.com/eaaCatalog/}{Catalog of Patterns of Enterprise Application Architecture}. The canonical enterprise application architecture reference, useful as the foundational pattern literature for comparing traditional enterprise patterns with agentic systems. \\
\end{tabularx}

\section{URL References}
\begin{itemize}[leftmargin=*]
    \item Salesforce: \url{https://architect.salesforce.com/docs/architect/fundamentals/guide/enterprise-agentic-architecture.html}
    \item Google Cloud: \url{https://docs.cloud.google.com/architecture/choose-design-pattern-agentic-ai-system}
    \item Martin Fowler: \url{https://martinfowler.com/eaaCatalog/}
    \item IBM: \url{https://www.ibm.com/think/topics/agentic-architecture}
    \item Google Cloud overview: \url{https://docs.cloud.google.com/architecture/agentic-ai-overview}
    \item Salesforce principles: \url{https://www.salesforce.com/news/stories/agentic-enterprise-design-principles/}
    \item Salesforce Architects homepage: \url{https://architect.salesforce.com}
\end{itemize}

\section{Why These Sources Matter}
Salesforce is especially useful because it explicitly frames agentic patterns as reusable architectural building blocks for enterprise systems, including agent layers, orchestration archetypes, security, observability, and lifecycle concerns. Google Cloud is valuable because it presents a clear pattern-selection guide with tradeoffs and practical guidance on when to apply or avoid certain designs. Martin Fowler's catalog provides the classic enterprise application architecture foundation, making it the best historical reference for the broader enterprise architecture context.

\section{Recommended Citation Set}
For a concise, authoritative bibliography, prioritize these three sources:
\begin{enumerate}[leftmargin=*]
    \item Salesforce, \textit{Enterprise Agentic Architecture and Design Patterns}.
    \item Google Cloud, \textit{Choose a design pattern for your agentic AI system}.
    \item Martin Fowler, \textit{Catalog of Patterns of Enterprise Application Architecture}.
\end{enumerate}

\section{Suggested Use}
This handout works well as:
\begin{itemize}[leftmargin=*]
    \item a reading list for enterprise AI architecture research,
    \item a handout for internal strategy or technical review meetings,
    \item a starting bibliography for a paper or slide deck on agentic enterprise systems.
\end{itemize}

\end{document}